\section{Simulation Setting}
\label{sec:setting}

This is a programming-based project, so what follows is the implementation and
measurement environment rather than a simulator.

\subsection{Programming environment}

The planner and cost model are Python, standard library only, so they run
anywhere and are testable without a GPU. The kernel is CUDA C++, compiled on
first use by PyTorch's \texttt{cpp\_extension} JIT and cached thereafter, and
driven from Python through a thin \texttt{profile()} / \texttt{apply()}
interface that uploads the cycle metadata once and then applies a plan
repeatedly.

The whole environment is pinned in a container built from the NGC PyTorch
image, which supplies torch, CUDA, \texttt{nvcc} and \texttt{ninja}. The
repository ships the \texttt{Dockerfile} and a dev container definition, so
the measurement environment is reproducible rather than described. Development
hardware is an RTX PRO 6000 Blackwell.

\subsection{Benchmark design}
\label{sec:benchmark}

A fair evaluation of tensor layout conversion has to cover a broad range of
tensor shapes rather than a small set of hand-picked examples. For a tensor of
$N$ elements and rank $r$, we define the shape space as every ordered
$r$-tuple whose dimensions multiply to $N$. Candidate shapes are enumerated by
prime-factorising $N$ and distributing the prime factors across the
dimensions:

\begin{equation}
s = \left(\prod_{i=1}^m p_i^{b_{i1}},\ \prod_{i=1}^m p_i^{b_{i2}},\ \ldots,\
\prod_{i=1}^m p_i^{b_{ir}}\right).
\label{eq:shapes}
\end{equation}

Here $N = p_1^{e_1} p_2^{e_2} \cdots p_m^{e_m}$, and $b_{ij}$ is the amount of
prime $p_i$ assigned to dimension $s_j$, so that
$\sum_{j} b_{ij} = e_i$ for every $i$. The number of ways to split one prime's
exponent across $r$ dimensions is given by stars and bars, so the size of the
space is
\begin{equation}
\Omega(N,r) = \prod_{i=1}^m \binom{e_i + r - 1}{\,r - 1\,}.
\label{eq:omega}
\end{equation}

This is the same construction the Cycle Splitter evaluation uses, which keeps
the two sets of results comparable. It also makes the comparison checkable.
The same equation builds the full list for any $N$ and any rank, so anyone can
generate the cases we did not pick and run them against any method here,
including ours.

We fix $N$ at \planElements{} elements and sweep ranks \planRankLo{} through
\planRankHi{}. From each rank's space we draw \planShapesPerRank{} shapes
uniformly, and for each shape we draw up to \planPermsPerShape{} non-identity
axis permutations. The identity is excluded because the tensor is already
contiguous under it. Shapes containing a unit axis are rejected, since an axis
of extent one can be moved without touching a byte. The space is large enough
that sampling is the only option. At rank \planRankHi{} it holds
\planSpaceTop{} shapes.

\subsection{Data sources}

\begin{description}
\item[Planner sweep.] Every permutation at each rank from \planRankLo{} to
  \planRankHi{}, planned by both strategies. CPU only, seeded, and committed
  to the repository, so every planner number in this proposal reproduces from
  a clean checkout with no GPU present. This is the source of
  Table~\ref{tab:planners} and Figure~\ref{fig:steps}.
\item[Throughput micro-benchmark.] A sweep of $D_{post}$ on the target GPU,
  used to locate the throughput knee $\tau$ and to fit $\mu$. It is
  regenerated per machine and not committed, because it describes the hardware
  that produced it.
\item[Pilot and benchmark sweeps.] Wider runs over rank-2 through rank-6
  shapes used to score the cost model against reality. These are hundreds of
  megabytes and are kept outside the repository, referenced through
  environment variables, with the scripts skipping rather than failing when
  they are absent.
\item[Comparison shapes.] For the head-to-head evaluation, the shapes reported
  by the closest prior systems~\cite{wu2025eithot,cheng2025ittpd}, so the
  comparison is on their ground rather than ours.
\end{description}

\subsection{Metrics}

\begin{description}
\item[Bytes moved.] The primary cost metric, and what the cost model
  predicts. Counted exactly from the cycle structure of each step, so it is
  independent of the machine and can be compared against a plan produced by
  any other strategy.
\item[Swap count.] Number of steps in the decomposition, which bounds the
  number of passes over the buffer.
\item[Auxiliary bytes.] Peak cycle metadata, which the objectives require to
  stay independent of $N$.
\item[Achieved throughput.] Bytes moved divided by measured wall time, in
  GB/s. This is what calibrates $\mu$ and what exposes the gap between the
  cost model's prediction and the hardware's behaviour.
\item[End-to-end latency and peak memory.] Measured against the out-of-place
  baseline, which is a permute followed by a contiguous copy. That is the
  operation this work is meant to replace. Latency says whether the approach
  is competitive. Peak memory is the claim it is actually making.
\end{description}

\subsection{Correctness}

Every plan is checked against the framework's own out-of-place permutation on
the same input. A decomposition that moves the fewest bytes is worthless if it
produces the wrong tensor, so correctness is a gate on every measurement
rather than a separate experiment.

\placeholder{State the shape and rank coverage of the correctness suite once
it is written, as a generated quantity rather than a typed number.}
